
\newpage
% El siguiente contenido corresponde a un segmento del manual de seguridad en radiaciones.
\section*{Falla en el mecanismo impulsor con la fuente en la manguera}

\subsection*{Clasificación}
% Según el documento existen diferentes clasificaciones dependiendo de la exposición de personas.
\begin{center}
\begin{tabular}{ | c | c | }
  \hline
No hay personas sobreexpuestas & I-A \\
  \hline
Hay personas sobreexpuestas & II-A\\
  \hline
\end{tabular}
\end{center}

\subsection*{Acciones recomendadas}
% Instrucciones paso a paso en caso de identificar el problema de falla en el mecanismo.
\begin{enumerate}
    \item \textbf{Acorde un área de 25 metros a la redonda de donde se encuentre la fuente:} Asegúrese de mantener la distancia para prevenir riesgos adicionales.
    \item \textbf{Evalúe la seguridad de las personas:} Observe si no hay riesgos. Si los hubiera, procure eliminarlos antes de pasar al siguiente punto.
    \item \textbf{Coloque el blindaje media caña sobre la manguera adelante del contenedor:} Esta medida es fundamental para proteger la integridad física de las personas y evitar la exposición directa a la fuente radiactiva.
\end{enumerate}

\begin{figure}[h]
\includegraphics[width=10cm]{mediaCana.png}
\centering
\end{figure}
